% \section{Commercial ReRAM and Timing Side-Channels}
% \label{sec:motivation}

% \subsection{Commercial ReRAM Deployment} Resistive random-access memory
% (ReRAM) refers to a class of non-volatile memory technologies that
% store data by modulating the electrical resistance of a switching
% layer. The most commercially mature variant is oxide-based ReRAM
% (OxRAM), which places a metal oxide switching layer between two
% electrodes and is CMOS compatible, enabling back-end-of-line
% integration in standard semiconductor processes~\cite{wong2012metal}.
% Unlike conventional flash memory, OxRAM offers low read latency,
% fine-grained addressability, and compatibility with standard CMOS
% processes~\cite{chenReRAMHistory2020}. These properties have driven its
% adoption as embedded data storage in low-power microcontrollers
% ~\cite{eenews2014panasonic_reram, nuvotonML23}. Additionally, ReRAM
% devices are intrinsically radiation tolerant, leading to significant
% study of ReRAM for use in aerospace applications. In the last decade,
% several commercial ReRAM products have entered mass production, making
% deployed devices widely accessible~\cite{fujreram4M, fujreram8M,
% rxreram12M, nuvotonML23}.

% Writing a value to a ReRAM cell involves growing or rupturing the CF
% within the switching layer through a sequence of voltage
% pulses~\cite{kawahara2013filament, yu2016resistive}. Filament formation
% and rupture processes are stochastic and sensitive to environmental
% conditions such as temperature and voltage
% fluctuations~\cite{ambrogio2013understanding,
% witzleben2017investigation, zhang2023deterministic}. Commercial
% architectures rely on a program-verify (PV) loop (and software
% polling/interupts) that iteratively applies programming pulses and
% checks cell resistance until the desired state is
% confirmed~\cite{kawahara2013filament, }. The duration of each PV loop
% is non-deterministic and depends on the initial physical state of the
% cells being programmed and the CF dynamics duing pulse applications.
% This cycle-to-cycle variability in write latency is an observable
% property of the device that carries information about its internal
% state.

% \subsection{Timing Properties and Security Primitives} The stochastic
% nature of ReRAM switching has attracted interest beyond data storage.
% Prior works have proposed exploiting write latency variability as an
% entropy source for hardware security primitives, including true random
% number generators (TRNGs) and physically unclonable functions
% (PUFs)~\cite{balatti2015true, jiang2017novel}. TRNGs provide the
% unpredictable seeds required by cryptographic protocols, and PUFs
% provide device-unique identifiers that resist cloning. The availability
% of commercial ReRAM in low-power embedded platforms makes it an
% attractive candidate for dual-use entropy harvesting, amortizing the
% cost of entropy generation across a component already serving a storage
% function. Commercial ReRAM devices implement error-correcting codes
% (ECCs) to improve reliability~\cite{hayakawa2017resolving,
% crafton2021cim}. Write operations in crossbar arrays are subject to
% strict sequencing constraints~\cite{xu2015overcoming}; the crossbar
% architecture requires SET and RESET operations to be serialized when
% they occur within the same word. ECC parity bits computed over user
% data introduce additional cell transitions on every write, determined
% by the codeword rather than the user-visible data alone. These
% transitions are invisible to the user but contribute to the write
% latency of every operation. The effect of ECC switching on write
% latency is not predictable without full knowledge of the underlying
% code. ECCs are considered vendor secrets, and their designers do not
% typically consider timing leakage as a design concern.

% \subsection{ReRAM Timing Leakage Channels} The stochastic PV loop and
% the deterministic ECC parity activity together produce structured write
% latency distributions that carry information about internal device
% state and embedded device architecture. These distributions constitute
% timing leakage channels: observable signals that correlate with
% information the device does not intentionally expose. Prior
% investigations of commercial ReRAM timing identified structured
% variance in these distributions but could not explain or model
% it~\cite{suri2018high, tolga2024trng}. The root cause, undocumented
% ECCs dominating the observable timing signal, was not identified. This
% gap has practical consequences. Without knowledge of the embedded ECC,
% empirical timing models are inaccurate, conflating deterministic parity
% activity with stochastic device physics. Security analyses of timing
% side-channels in crossbar memories have overlooked the dominant
% contributor to write latency variance~\cite{khan2021comprehensive,
% kommareddy2019crossbar, wang2023side, krautter2022data}. Entropy claims
% for ReRAM-based TRNGs are overstated, as structured ECC perturbations
% are mistaken for true device stochasticity. The timing side-channel
% attack surface of commercial ReRAM products cannot be fully
% characterized without first recovering these hidden architectural
% features.

% \subsection{Research Gaps} \label{sec:research_gaps_and_motivation} The
% only existing technique for bit-exact ECC recovery from a commercial
% memory product is BEER~\cite{patel2020bit}, which targets DRAM by
% exploiting data retention failures that expose ECC corrections through
% direct readback. Commercial ReRAM products do not support this
% mechanism; the PV loop ensures every write converges to the target
% state, preventing the retention-based error injection that BEER relies
% upon, and parity bits are not directly readable. Physical reverse
% engineering of commercial ReRAM has been demonstrated using
% conductive-probe atomic force microscopy~\cite{tay2023afm}, but this
% approach is destructive, costly, and impractical outside a laboratory
% setting. No non-invasive technique for recovering the ECC of a
% commercial ReRAM product exists. With knowledge of the underlying ECC,
% accurate codeword-level timing models become realizable, the stochastic
% component of write latency can be isolated and rigorously characterized
% as an entropy source, and the timing side-channel attack surface of
% commercial ReRAM products can be fully analyzed for the first time.

\paragraph{Thesis Statement}
ReRAM defect dynamics share many similarities with those exploited by
FPGA pentimenti attacks. ReRAM defect dynamics leak through write
latency via a stochastic timing component. However, commercial ReRAM
write latency decomposes into structural and stochastic components, but
both are recoverable through non-invasive timing analysis. This
decomposition enables accurate timing models, rigorous entropy
characterization, and the measurement capabilities required for study of
pentimenti-class side-channel analysis in commercial ReRAM modules.
